\par\noindent\textbf{(1)} 外力矩为零，轴向角动量 \(L=I(t)\omega(t)\) 为常数，因此 \(\omega(t)=L/I(t)\)。微分给 \(\dot\omega=-(\dot I/I)\omega\)。尽管没有外扭矩，角速度仍可随着内部质量径向运动而改变。
\par\noindent\textbf{(2)} 转动能量 \(K=L^2/(2I)\)。沿过程微分 \(\dot K=-L^2\dot I/(2I^2)=-\tfrac12\omega^2\dot I\)。若质量向内收缩 \(\dot I<0\)，旋转能量上升，来源是内部执行机构对滑块做正功而非无中生有。
\par\noindent\textbf{(3)} 若轴上本体惯量 \(I_0\)，另有两个质量为 \(m\) 的滑块位于半径 \(r\)，则 \(I(r)=I_0+2mr^2\)。从 \(r_i\) 移至 \(r_f\)，角速度比 \(\omega_f/\omega_i=(I_0+2mr_i^2)/(I_0+2mr_f^2)\)，可与角动量直接检验。
\par\noindent\textbf{(4)} 在固定总角动量下，\(dK/dr=-L^2(4mr)/(2I^2)=-2mr\omega^2\)。向内慢移需抵抗离心效应，机构的功 \(W=K_f-K_i=\int_{r_i}^{r_f}(-2mr\omega^2)dr\)。若过程有滑动摩擦，还必须加入耗散热，不能只用转动能账。
